\section{Глава 3. Конечные марковские процессы принятия решений}
\subsection{Упражнение 3.1(c. 78)}
Придумайте три примера задач, укладывающихся в формализм
МППР. Определите для каждого состояния, действия и  вознаграждения. Постарайтесь, чтобы примеры как можно сильнее различались. Формализм абстрактный и гибкий, он может применяться разными способами. Хотя бы в одном примере попытайтесь подобраться к его пределам.
\begin{solution}
    \textbf{Первый пример.} Начнем с примитивного. Рассмотрим популярную в наше время задачу беспилотных автомобилей. Действиями могут быть следующие параметры: угол, на который следует повернуть руль, сила нажатия на педали тормоза/газа. Состоянием могут быть как характеристики движения, например, скорость, так и сцена возле автомобиля, полученная с лидара. Для вознаграждений же можно придумать множество реализаций. Например, это может быть сложной функцией от расстояния до ближайшего объекта. Разумно будет добавить штраф за отклонение с полосы или резкие ускорения. \\
    \textbf{Второй пример.} Выбор партнера для отношений. Поскольку у нас нет ограничений на временные шаги, то можем рассматривать тройку (состояние, действие, награда) с частотой, например, в месяц. Приведем возможные примеры действий: подарить цветы, пригласить на свидание, оплатить совместный отдых или расстаться и сменить партнера. Состояние включает в себя характеристики партнера(психические, физические), накопленные обиды, вашу удовлетворенность, ваш социальный/карьерный рост за время этих отношений. А награду введем +1000000 за заключение брачного договора, +1000 за преодоление трудностей, -1000000000000 за разбитые сердца, +100 за стабильность в отношениях. Это пример нестационарной среды, в которой распределение вероятностей зависит от текущего партнера.\\
    \textbf{Третий пример.} Рассмотрим задачу эволюции целой цивилизации как единого агента, принимающего решения раз в столетие. Действиями являются выбор научных парадигм, политических устройств, моральных норм, направлений экспансии или отказа от неё. Состояние включает суммарное знание, технологический уровень, численность населения, распределение ресурсов, устойчивость экосистемы и культурные ценности. Награда определяется выживанием цивилизации, ростом совокупного благополучия, расширением доступной Вселенной или сохранением сознательной жизни, а терминальным состоянием может быть вымирание, технологическая сингулярность или тепловая смерть локального региона космоса.
\end{solution}

\subsection{Упражнение 3.2(c. 79)}
Достаточно ли формализма МППР для полезного представления
всех задач обучения целеустремленного агента? Можете ли вы придумать очевидные исключения? 
\begin{solution}
    Во-первых, далеко не во всех задачах имеет место марковское свойство. Контрпример - покер. В этой игре $\mathbb{P}(S_{t+1}, R_{t+1} \mid S_t, A_t)$ зависит и от выбора соперника, и от прошлой истории.\\
    Во-вторых, в классической формулировке МППР подразумевается стационарность. Если среда нестационарна(те же самообучающиеся крестики-нолики из Главы $1$), то мы имеем дело с $\mathbb{P}_t(S_{t+1}, R_{t+1} \mid S_t, A_t)$. \\
    В-третьих, существует огромный круг задач, где использование награды не самым корректным образом отражает структуру задачи. Например, если рассматривать постановки, в которых человек является агентом(как в предыдущем упражнении), то его цель может не соответствовать любой предложенной награде.
\end{solution}

\subsection{Упражнение 3.3(c. 79)}
Рассмотрим задачу о  вождении автомобиля. Можно было бы
определить действия в терминах акселератора, руля и тормозов, т.е. тех мест, где ваше тело входит в контакт с машиной. А можно было бы поместить их дальше снаружи – например, там, где шина контактирует с дорогой, считая вашими действиями крутящий момент на колесе. Или поместить их дальше внутри – скажем, там, где ваш мозг вступает в контакт с телом, и считать действиями сокращения
мускулов, управляющих вашими конечностями. Или можно перейти на совсем уж высокий уровень и сказать, что ваше действие – это выбор места назначения. Какой уровень правильный, т. е. где нужно было бы провести границу между агентом и окружающей средой? На каком основании следует предпочесть одно положение границы другому? Существует ли фундаментальная причина предпочесть одно положение другому или это свободный выбор?
\begin{solution}
    Граница между агентом и средой в задаче вождения автомобиля не имеет единственного "правильного"\ положения и определяется не физикой, а целями моделирования. Предпочтение одного уровня другому носит прагматический характер: выбранная граница должна обеспечивать марковость состояния, соответствовать реально управляемым агентом переменным и приводить к разумной сложности динамики и обучения. Таким образом, фундаментальной причины предпочесть один уровень другому нет - это свободный, но осмысленный моделировочный выбор, зависящий от целей и удобства анализа.
\end{solution}

\subsection{Упражнение 3.4(c. 80)}
Постройте такую же таблицу, как в примере 3.3, но для $p(s', r \mid s, a)$. В ней должны быть столбцы $s, a, s', r$ и $p(s', r \mid s, a)$ и по одной строке для каждой четверки такой, что $p(s', r \mid s, a) > 0$.
\begin{solution}
\leavevmode\par
\medskip
    \begin{center}
    \begin{tabular}{c c c c c}
    \hline
    $s$ & $a$ & $s'$ & $r$ & $p(s',r\mid s,a)$ \\
    \hline
    high & search & high & $r_{\text{search}}$ & $\alpha$ \\
    high & search & low  & $r_{\text{search}}$ & $1-\alpha$ \\
    \hline
    high & wait   & high & $r_{\text{wait}}$   & $1$ \\
    \hline
    low  & search & low  & $r_{\text{search}}$ & $\beta$ \\
    low  & search & high & $-3$                & $1-\beta$ \\
    \hline
    low  & wait   & low  & $r_{\text{wait}}$   & $1$ \\
    \hline
    low  & recharge & high & $0$               & $1$ \\
    \hline
    \end{tabular}
    \end{center}
\end{solution}

\subsection{Упражнение 3.5(c. 83)}
Равенства в разделе $3.1$ справедливы для непрерывного случая, а для эпизодических задач их необходимо модифицировать (совсем чуть-чуть). Покажите, что вы знаете, как именно, для чего приведите модифицированный вариант формулы $(3.3)$. 
\begin{equation}
\sum_{s' \in S} \sum_{r \in R}
p(s', r \mid s, a) = 1,
\qquad
\forall\, s \in S,
\ \forall\, a \in A(s).
\tag{3.3}
\end{equation}
\begin{solution}
    \[
    \sum_{s' \in S^+} \sum_{r \in R}
    p(s', r \mid s, a) = 1,
    \qquad
    \forall\, s \in S,
    \ \forall\, a \in A(s).
    \]
\end{solution}

\subsection{Упражнение 3.6(c. 84)}
Предположим, что балансирование стержня рассматривается как эпизодическая задача, но с  использованием обесценивания, так что все вознаграждения равны $0$, кроме вознаграждения за неудачу, равного $-1$. Каким тогда будет доход в каждый момент времени? Чем этот доход отличается от дохода в непрерывной задаче с обесцениванием?
\begin{solution}
    В такой постановке задачи доход также будет определяться величиной $-\gamma^K$, где $K$ - количество временных шагов до наступления неудачи. В случае эпизодической задачи 
    \[G_t = -\gamma^K\]
    После неудачи эпизод заканчивается и начинается следующий. В непрерывной же задаче переключения эпизода не происходит и доход выглядит следующим образом:
    \[G_t = -\gamma^{K_1} - \gamma^{K_2} - \dots - \gamma^{K_n} - \dots,\]
    где $K_i$ - момент неудачи в $i$-ой попытке.
\end{solution}

\subsection{Упражнение 3.7(c. 84)}
Допустим, вы проектируете робота для прохождения лабиринта. Вы решили давать вознаграждение $+1$ за выход из лабиринта и $0$ в остальных случаях. Похоже, задача естественно распадается на эпизоды – последовательные прохождения лабиринта, поэтому вы решили рассматривать ее как эпизодическую, в которой цель – максимизировать ожидаемое полное вознаграждение $(3.7)$. Погоняв некоторое время обучающегося агента, вы обнаружили полное отсутствие прогресса. Что не так? Достаточно ли доходчиво вы сообщили агенту, чего именно он должен достичь?
\begin{solution}
    Запишем формулу дохода для случайного эпизода:
    \[G_t = r_{t+1} + r_{t+2} + \dots + r_T\]
    Вышеупомянутое определение вознаграждения вне зависимости от длины эпизода $T$ приводит к тому, что
    \[r_{t+1} = r_{t+2} = \dots = r_{T-1} = 0 \quad \text{и} \quad r_T = 1\]
    Тогда в любом эпизоде $G_t = 1$. Из этого следует, что награда никак не согласуется со временем, потраченным на поиск выхода из лабиринта. С точки зрения вознаграждения эпизод, когда робот долго не мог найти выход и блуждал по лабиринту, будет равен эпизоду, когда робот нашел выход за минимальное число шагов. Решением будет, например, добавить дисконтирование($\gamma$).
\end{solution}

\subsection{Упражнение 3.8(c. 84)}
Предположим, что $\gamma = 0.5$, последовательность полученных вознаграждений имеет вид: $R_1 = -1, R_2 = 2, R_3 = 6, R_4 = 3, R_5 = 2 \ \text{и} \ T = 5$. Чему равны $G_0, G_1, \dots , G_5$? Указание: идите от конца к началу.
\begin{solution}
    Используем рекурсивное определение дохода:
    \[
    G_t = R_{t+1} + \gamma G_{t+1}, 
    \qquad G_T = 0.
    \]
    
    Дано: 
    \[
    \gamma = 0.5, \quad 
    R_1=-1,\; R_2=2,\; R_3=6,\; R_4=3,\; R_5=2, \quad T=5.
    \]
    
    Вычисляем с конца.
    
    \[
    G_5 = 0.
    \]
    
    \[
    G_4 = R_5 + \gamma G_5
         = 2 + 0.5 \cdot 0
         = 2.
    \]
    
    \[
    G_3 = R_4 + \gamma G_4
         = 3 + 0.5 \cdot 2
         = 3 + 1
         = 4.
    \]
    
    \[
    G_2 = R_3 + \gamma G_3
         = 6 + 0.5 \cdot 4
         = 6 + 2
         = 8.
    \]
    
    \[
    G_1 = R_2 + \gamma G_2
         = 2 + 0.5 \cdot 8
         = 2 + 4
         = 6.
    \]
    
    \[
    G_0 = R_1 + \gamma G_1
         = -1 + 0.5 \cdot 6
         = -1 + 3
         = 2.
    \]
    
    Итак,
    \[
    G_0 = 2, \quad
    G_1 = 6, \quad
    G_2 = 8, \quad
    G_3 = 4, \quad
    G_4 = 2, \quad
    G_5 = 0.
    \]
\end{solution}

\subsection{Упражнение 3.9(c. 84)}
Предположим, что $\gamma = 0.9$ и последовательность вознаграждений начинается с $R_1 = 2$, за которым следует бесконечная последовательность чисел $7$.
Чему равны $G_0$ и $G_1$? 
\begin{solution}
    По определению дохода:
    \[
    G_t = \sum_{k=0}^{\infty} \gamma^k R_{t+k+1}.
    \]
    Тогда
    \[
    G_1 = 7 + 0.9 \cdot 7 + 0.9^2 \cdot 7 + \dots = 
    7 \cdot \sum_{k=0}^{\infty} (0.9)^k
     = 7 \cdot \frac{1}{1-0.9}
     = 7 \cdot \frac{1}{0.1}
     = 70.
    \]
    Пользуясь рекурсивной формулой,
    \[G_0 = R_1 + \gamma \cdot G_1 
     = 2 + 0.9 \cdot 70
     = 2 + 63
     = 65.\]
\end{solution}

\subsection{Упражнение 3.10(c. 84)}
Докажите равенство $(3.10)$.
\begin{equation}
G_t = \sum_{k=0}^{\infty} \gamma^k = \frac{1}{1 - \gamma}, \qquad 0 < \gamma < 1
\tag{3.10}
\end{equation}
\begin{solution}
    Рассмотрим частичную сумму геометрического ряда:
    \[
    S_n=\sum_{k=0}^{n}\gamma^k.
    \]
    Из формулы суммы геометрической прогрессии:
    \[
    S_n=\frac{1-\gamma^{\,n+1}}{1-\gamma}.
    \]
    Так как $\gamma<1$, то $\gamma^{\,n+1}\to 0$ при $n\to\infty$. Следовательно,
    \[
    \sum_{k=0}^{\infty}\gamma^k
    =\lim_{n\to\infty} S_n
    =\lim_{n\to\infty}\frac{1-\gamma^{\,n+1}}{1-\gamma}
    =\frac{1}{1-\gamma}.
    \]
\end{solution}

\subsection{Упражнение 3.11(c. 86)}
Если текущее состояние равно $S_t$ и действия выбираются согласно стохастической стратегии $\pi$, то каково математическое ожидание $R_{t+1}$
в терминах $\pi$ и функции $p$ с четырьмя аргументами (3.2)?
\begin{equation}
p(s', r \mid s, a)
=
\mathbb{P}\!\bigl\{ S_t = s',\, R_t = r \,\big|\, S_{t-1} = s,\, A_{t-1} = a \bigr\}.
\tag{3.2}
\end{equation}
\begin{solution}
    \[
    \begin{aligned}
    \mathbb{E}\!\left[R_{t+1}\mid S_t=s\right]
    &= \sum_{r\in R} r\,\mathbb{P}\!\left(R_{t+1}=r\mid S_t=s\right) \\
    &= \sum_{r\in R} r \sum_{a\in A(s)} \pi(a\mid s)
    \sum_{s'\in S} p(s',r \mid s,a).
    \end{aligned}
    \]
\end{solution}

\subsection{Упражнение 3.12(c. 86)}
Выразите $v_\pi$ через $q_\pi$ и $\pi$.
\[
    v_{\pi}(s)
    =
    \mathbb{E}_{\pi}\!\left[ G_t \mid S_t = s \right]
    =
    \mathbb{E}_{\pi}\!\left[ \left. \sum_{k=0}^{\infty} \gamma^k R_{t+k+1} \,\right|\, S_t = s \right],
    \quad \text{для всех } s \in S.
\]
\[
    q_{\pi}(s,a)
    =
    \mathbb{E}_{\pi}\!\left[ G_t \mid S_t = s, A_t = a \right]
    =
    \mathbb{E}_{\pi}\!\left[ \left. \sum_{k=0}^{\infty} \gamma^k R_{t+k+1}
    \,\right|\, S_t = s, A_t = a \right].
\]
\begin{solution}
    \begin{align*}
    v_\pi(s)
    &= \mathbb{E}_\pi\!\left[ G_t \mid S_t = s \right] =\\
    &= \sum_{a\in A(s)} \mathbb{P}_\pi(A_t=a \mid S_t=s)\,
       \mathbb{E}_\pi\!\left[ G_t \mid S_t=s, A_t=a \right] =\\
    &= \sum_{a\in A(s)} \pi(a\mid s)\, q_\pi(s,a).
    \end{align*}
\end{solution}

\subsection{Упражнение 3.13(c. 86)}
Выразите $q_\pi$ через $v_\pi$ и функцию $p$ с четырьмя аргументами. .
\[
    v_{\pi}(s)
    =
    \mathbb{E}_{\pi}\!\left[ G_t \mid S_t = s \right]
    =
    \mathbb{E}_{\pi}\!\left[ \left. \sum_{k=0}^{\infty} \gamma^k R_{t+k+1} \,\right|\, S_t = s \right],
    \quad \text{для всех } s \in S.
\]
\[
    q_{\pi}(s,a)
    =
    \mathbb{E}_{\pi}\!\left[ G_t \mid S_t = s, A_t = a \right]
    =
    \mathbb{E}_{\pi}\!\left[ \left. \sum_{k=0}^{\infty} \gamma^k R_{t+k+1}
    \,\right|\, S_t = s, A_t = a \right].
\]
\[
p(s', r \mid s, a)
=
\mathbb{P}\!\bigl\{ S_t = s',\, R_t = r \,\big|\, S_{t-1} = s,\, A_{t-1} = a \bigr\}.
\]
\begin{solution}
    \[
    q_\pi(s,a)
    = \mathbb{E}_\pi\!\left[ G_t \mid S_t=s, A_t=a \right]
    = \sum_{s'\in S}\sum_{r\in R} p(s',r\mid s,a)\,\Bigl(r + \gamma\, v_\pi(s')\Bigr).
    \]
\end{solution}

\subsection{Упражнение 3.14(c. 89)}
Уравнение Беллмана $(3.14)$ должно выполняться в  каждом состоянии для функции ценности $v_{\pi}$, показанной на рис. $3.2$ (справа) в примере  $3.5$. Покажите численно, что это уравнение выполняется для центрального состояния с ценностью $+0.7$ по отношению к четырем его соседним состояниям
с ценностями $+2.3$, $+0.4$, $-0.4$ и $+0.7$. (Точность этих чисел – всего один знак после запятой.)
\begin{equation}
    v_\pi(s)
    =
    \sum_{a} \pi(a \mid s)
    \sum_{s',r} p(s', r \mid s, a)
    \bigl[ r + \gamma v_\pi(s') \bigr]
\tag{3.14}
\end{equation}
\begin{solution}
    Для примера 3.5 стратегия равновероятная, поэтому
    \[
    \pi(a\mid s)=\frac14,\qquad a\in\{\text{север, юг, запад, восток}\},
    \]
    а для центрального состояния все переходы детерминированы в соседние клетки и дают награду \(r=0\).
    Также в примере используется \(\gamma=0.9\).
    
    Тогда уравнение Беллмана (3.14) для этого состояния принимает вид
    \[
    v_\pi(s)=\frac14 \sum_{a}\bigl[0+\gamma v_\pi(s_a)\bigr]
    =\frac{\gamma}{4}\Bigl(v_\pi(s_{\uparrow})+v_\pi(s_{\downarrow})+v_\pi(s_{\leftarrow})+v_\pi(s_{\rightarrow})\Bigr).
    \]
    
    Подставим значения соседей \(2.3,\;0.4,\;-0.4,\;0.7\):
    \[
    v_\pi(s)=\frac{0.9}{4}\,(2.3+0.4-0.4+0.7)
    =\frac{0.9}{4}\cdot 3.0
    =0.675 \approx 0.7.
    \]
    
    Получили \(v_\pi(s)\approx 0.7\), что согласуется с указанной на рисунке величиной (с точностью до одного знака после запятой).
\end{solution}

\subsection{Упражнение 3.15(c. 89)}
В примере сеточного мира вознаграждения положительны для
целей, отрицательны для выхода на края мира и равны нулю в остальных случаях. Что важно: знаки вознаграждений или только расстояния между ними? Докажите, пользуясь формулой $(3.8)$, что прибавление константы c ко всем вознаграждениям прибавляет константу $v_c$ к ценностям всех состояний и потому не влияет на относительные ценности состояний при любой стратегии. Выразите $v_c$
 через $c$ и $\gamma$.
\begin{equation}
    G_t = R_{t+1} + \gamma R_{t+2} + \gamma^2 R_{t+3} + \cdots
= \sum_{k=0}^{\infty} \gamma^k R_{t+k+1}.
\tag{3.8}
\end{equation}
\begin{solution}
    Рассмотрим
    \[
    R'_{t+k+1}=R_{t+k+1}+c,\qquad k=0,1,2,\dots
    \]
    Тогда по формуле (3.8) новый доход
    \[
    \begin{aligned}
    G'_t
    &= \sum_{k=0}^{\infty}\gamma^k R'_{t+k+1}
    = \sum_{k=0}^{\infty}\gamma^k (R_{t+k+1}+c) \\
    &= \sum_{k=0}^{\infty}\gamma^k R_{t+k+1}
    + c\sum_{k=0}^{\infty}\gamma^k
    = G_t + c\sum_{k=0}^{\infty}\gamma^k.
    \end{aligned}
    \]
    При $0 < \gamma<1$ имеем геометрическую сумму
    \[
    \sum_{k=0}^{\infty}\gamma^k=\frac{1}{1-\gamma},
    \]
    поэтому
    \[
    G'_t=G_t+\frac{c}{1-\gamma}.
    \]
    Обозначим ценность состояния:
    \[
    v_\pi(s)=\mathbb{E}_\pi[G_t\mid S_t=s],\qquad
    v'_\pi(s)=\mathbb{E}_\pi[G'_t\mid S_t=s].
    \]
    Берем математическое ожидание:
    \[
    v'_\pi(s)=\mathbb{E}_\pi\!\left[G_t+\frac{c}{1-\gamma}\,\middle|\,S_t=s\right]
    =\mathbb{E}_\pi[G_t\mid S_t=s]+\frac{c}{1-\gamma}
    =v_\pi(s)+v_c,
    \]
    где
    \[
    v_c=\frac{c}{1-\gamma}.
    \]
    Следовательно, для любых $s_1,s_2$:
    \[
    v'_\pi(s_1)-v'_\pi(s_2)=(v_\pi(s_1)+v_c)-(v_\pi(s_2)+v_c)=v_\pi(s_1)-v_\pi(s_2),
    \]
    то есть прибавление константы $c$ ко всем вознаграждениям сдвигает значения всех состояний на одну и ту же константу $v_c$
    и не влияет на их относительные ценности при любой стратегии $\pi$.
\end{solution}

\subsection{Упражнение 3.16(c. 89)}
Теперь рассмотрим прибавление константы $c$ ко всем вознаграждениям в эпизодической задаче, например о прохождении лабиринта. Будет ли это иметь какой-то эффект или задача останется такой же, как в случае непрерывной задачи выше? Почему? Приведите пример. 
\begin{solution}
Рассмотрим обесцененную версию эпизодической задачи. Пусть эпизод завершается за случайное число шагов $T$.
Исходный доход:
\[
G_t=\sum_{k=0}^{T-t-1}\gamma^k R_{t+k+1}.
\]
После сдвига вознаграждений $R'_{t+k+1}=R_{t+k+1}+c$ получаем
\[
G'_t
=\sum_{k=0}^{T-t-1}\gamma^k (R_{t+k+1}+c)
=G_t + c\sum_{k=0}^{T-t-1}\gamma^k.
\]
Геометрическая сумма теперь зависит от длины эпизода:
\[
\sum_{k=0}^{T-t-1}\gamma^k
=\frac{1-\gamma^{T-t}}{1-\gamma}.
\]
Следовательно,
\[
G'_t
=G_t + \frac{c(1-\gamma^{T-t})}{1-\gamma}.
\]
В отличие от непрерывной задачи, добавка зависит от $T-t$,
то есть от того, сколько шагов осталось до завершения.
Поэтому разные траектории получают разные добавки,
и относительные ценности состояний меняются.

\textbf{Пример.}
Рассмотрим задачу лабиринта с $\gamma=1$, нулевыми наградами на каждом шагеи $+1$ в терминальном состоянии.
Тогда $G_t=1$ независимо от длины пути.

Теперь прибавим $c=1$ ко всем наградам. Тогда каждый шаг даёт $+1$, а в конце ещё $+1$.
Доход равен:
\[
G'_t = T-t+1.
\]
Теперь выгодно затягивать эпизод как можно дольше.
Оптимальная стратегия изменилась.

Следовательно, в эпизодических задачах прибавление константы ко всем вознаграждениям меняет задачу, поскольку дополнительная сумма зависит от длины эпизода.
\end{solution}

\subsection{Упражнение 3.17(c. 90)}
\begin{wrapfigure}{r}{0.3\textwidth}
    \centering
    \includegraphics[width=0.25\textwidth]{q_pi_graph.png}
\end{wrapfigure}
Как выглядит уравнение Беллмана для ценностей действий, т. е. для $q_{\pi}$? Оно должно выражать ценность действия $q_{\pi}(s, a)$ в терминах ценностей действий $q_{\pi}(s', a')$ возможных преемников пары состояние–действие $(s, a)$. Указание: этому уравнению соответствует диаграмма предшествующих состояний справа. Выпишите последовательность равенств, аналогичную $(3.14)$, но для ценностей действий.
\begin{solution}
    Воспользуемся результатом упражнений 3.12 и 3.13:
    \begin{align*}
    q_\pi(s,a)
    &= \mathbb{E}_\pi\!\left[ G_t \mid S_t=s, A_t=a \right] = \\
    &= \sum_{s'\in S}\sum_{r\in R} p(s',r\mid s,a)\,\Bigl(r + \gamma\, v_\pi(s')\Bigr)  = \\
    &= \sum_{s'\in S}\sum_{r\in R} p(s',r\mid s,a)\,\Bigl(r + \gamma\sum_{a'\in A(s')} \pi(a'\mid s')\, q_\pi(s',a')\Bigr) 
    \end{align*}
\end{solution}

\subsection{Упражнение 3.18(c. 90)}
Ценность состояния зависит от ценностей действий, возможных в этом состоянии, и от вероятностей выбора
действий при текущей стратегии. Мы можем представить это в виде небольшой диаграммы предшествующих состояний с корнем в состоянии, на которой показаны все возможные действия:
\begin{figure}[H]
    \centering
    \includegraphics[width=0.8\linewidth]{fig_3_18_task.png}
\end{figure}
Выпишите уравнение, соответствующее этому интуитивному представлению и диаграмме для ценности в корневом узле $v_{\pi}(s)$, в терминах ценности в ожидаемом
листовом узле $q_{\pi}(s, a)$, при условии что $S_t = s$. Это уравнение должно включать математическое ожидание, обусловленное следованием стратегии $\pi$. Затем выпишите
второе уравнение, в котором ожидаемая ценность выражена явно через $\pi(a | s)$, так чтобы никакого математического ожидания ценности в уравнение не входило. 
\begin{solution}
\[
v_\pi(s)=\mathbb{E}_\pi\!\left[q_\pi(s,A_t)\mid S_t=s\right].
\]
Если множество действий дискретно, то это ожидание можно выписать явно через вероятности выбора действий:
\[
v_\pi(s)=\sum_{a}\pi(a\mid s)\,q_\pi(s,a).
\]
\end{solution}

\subsection{Упражнение 3.19(c. 91)}
Ценность действия $q_{\pi}(s, a)$ зависит от следующего ожидаемого вознаграждения и математического ожидания суммы оставшихся вознаграждений. И снова мы можем рассуждать о задаче, опираясь на небольшую диаграмму предшествующих состояний, на этот раз с корнем в действии (паре 
состояние-действие), которое ведет в одно из следующих возможных состояний:
\begin{figure}[H]
    \centering
    \includegraphics[width=0.7\linewidth]{fig_3_19_task.png}
\end{figure}
Выпишите уравнение, соответствующее этому интуитивному представлениюи диаграмме для ценности действия 
$q_{\pi}(s, a)$, в терминах следующего ожидаемого
вознаграждения $R_{t+1}$ и ожидаемой ценности следующего состояния $v_{\pi}(S_{t+1})$, при условии что $S_t = s$ и $A_t = a$. Это уравнение должно включать математическое
ожидание, но \emph{не} обусловленное следованием стратегии. Затем выпишите второе уравнение, в  котором ожидаемая ценность выражена явно через величину $p(s', r \mid s, a)$, определенную в $(3.2)$, так чтобы никакого математического ожидания ценности в уравнение не входило.
\begin{equation}
p(s', r \mid s, a)
=
\mathbb{P}\!\bigl\{ S_t = s',\, R_t = r \,\big|\, S_{t-1} = s,\, A_{t-1} = a \bigr\}.
\tag{3.2}
\end{equation}
\begin{solution}
При фиксированных $S_t=s$ и $A_t=a$ случайность обусловлена только динамикой среды, поэтому
\[
q_\pi(s,a)=\mathbb{E}\!\left[\,R_{t+1}+\gamma\,v_\pi(S_{t+1}) \,\middle|\, S_t=s,\ A_t=a\right].
\]
Расписывая это ожидание явно через модель среды $p(s',r\mid s,a)$ из (3.2), получаем
\[
q_\pi(s,a)=\sum_{s'}\sum_{r} p(s',r\mid s,a)\,\Bigl(r+\gamma\,v_\pi(s')\Bigr).
\]
\end{solution}

\subsection{Упражнение 3.20(c. 96)}
Нарисуйте или опишите оптимальную функцию ценности состояний для примера игры в гольф. 
\begin{solution}
    В примере 3.7 описана оптимальная стратегия: два удара драйвером и один паттером. Оптимальная функция ценности состояний соответствует именно этой политике. Поэтому 
    \[
    v_*(s) =
    \begin{cases}
    -1, & \text{в области грина}, \\
    -2, & \text{где $q_*(s, \text{driver})=-2$}, \\
    -3, & \text{где $q_*(s, \text{driver})=-3$}, \\
    \dots
    \end{cases}
    \]
\end{solution}

\subsection{Упражнение 3.21(c. 96)}
Нарисуйте или опишите линии уровня оптимальной функции
ценности действий для игры паттером, $q_*(s,\text{putter})$, в примере игры в гольф.
\begin{solution}
    Не ограничивая общности, предположим, что линии уровня $-2, -3$ функции $q_*(s, \text{driver})$ совпадают  с линиями уровня $-3, -6$ функции $q_*(s, \text{putter})$ соответственно. \\
    Выбор действия согласно $q_*(s,\text{putter})$ означает, что из состояния $s$ первый шаг мы делаем паттером, а дальше действуем оптимальным образом. \\
    Поэтому 
    \[
    q_*(s,\text{putter}) =
    \begin{cases}
    -1, & \text{в области грина}, \\
    -2, & \text{где $v_{\text{putt}}(s)=-2$}, \\
    -3, & \text{где $-4 \le v_{\text{putt}}(s) \le -3$}, \\
    -4, & \text{где $-7 \le v_{\text{putt}}(s) \le -5$}, \\
    \dots
    \end{cases}
    \]
\end{solution}

\subsection{Упражнение 3.22(c. 96)}
\begin{wrapfigure}{r}{0.3\textwidth}
    \centering
    \includegraphics[width=0.25\textwidth]{fig_3_22_task.png}
\end{wrapfigure}
Рассмотрим непрерывный МППР, показанный на рисунке справа. Единственное решение принимается в верхнем состоянии, в котором возможно два действия, $left$ и $right$. Числа обозначают детерминированные вознаграждения за каждое действие. Имеется ровно две детерминированные стратегии, $\pi_{left}$ и $\pi_{right}$. Какая стратегия оптимальна, если $\gamma = 0$? А если $\gamma = 0.9$? А если $\gamma = 0.5$?
\begin{solution}
    Решение принимается только в состоянии $s_0$. После выбора действия процесс детерминированно возвращается в $s_0$ за два шага.

    \textbf{Стратегия $\pi_{left}$.}
    
    Последовательность наград:
    \[
    (+1,\,0,\,+1,\,0,\,+1,\,0,\dots)
    \]
    
    Тогда
    \[
    v_{left}(s_0)
    =
    1 + \gamma^2 + \gamma^4 + \dots
    =
    \frac{1}{1-\gamma^2}.
    \]
    
    \textbf{Стратегия $\pi_{right}$.}
    
    Последовательность наград:
    \[
    (0,\,+2,\,0,\,+2,\,0,\,+2,\dots)
    \]
    
    Тогда
    \[
    v_{right}(s_0)
    =
    2\gamma + 2\gamma^3 + 2\gamma^5 + \dots
    =
    \frac{2\gamma}{1-\gamma^2}.
    \]
    
    \textbf{Подсчёт для заданных значений $\gamma$:}
    
    \[
    \gamma = 0:
    \quad
    v_{left} = 1,
    \quad
    v_{right} = 0.
    \]
    
    \[
    \gamma = 0.9:
    \quad
    v_{left} = \frac{1}{1-0.81} = \frac{1}{0.19} \approx 5.26,
    \quad
    v_{right} = \frac{1.8}{0.19} \approx 9.47.
    \]
    
    \[
    \gamma = 0.5:
    \quad
    v_{left} = \frac{1}{1-0.25} = \frac{1}{0.75} = \frac{4}{3},
    \quad
    v_{right} = \frac{1}{0.75} = \frac{4}{3}.
    \]
\end{solution}

\subsection{Упражнение 3.23(c. 96)}
Выпишите уравнение Беллмана для $q_*$ в задаче о роботе-уборщике.
\begin{solution}
    Уравнение оптимальности Беллмана:
    \[
    q_*(s,a)=\sum_{s', r} p(s',r\mid s,a)\Bigl[r(s,a,s')+\gamma \max_{a'} q_*(s',a')\Bigr].
    \]
    
    Для робота-уборщика:
    \[
    q_*(h,s)
    =\alpha\Bigl[r_s+\gamma \max_{a'} q_*(h,a')\Bigr]
    +(1-\alpha)\Bigl[r_s+\gamma \max_{a'} q_*(l,a')\Bigr].
    \]
    
    \[
    q_*(h,w)
    =r_w+\gamma \max_{a'} q_*(h,a').
    \]
    
    \[
    q_*(l,s)
    =\beta\Bigl[r_s+\gamma \max_{a'} q_*(l,a')\Bigr]
    +(1-\beta)\Bigl[-3+\gamma \max_{a'} q_*(h,a')\Bigr].
    \]
    
    \[
    q_*(l,w)
    =r_w+\gamma \max_{a'} q_*(l,a').
    \]
    
    \[
    q_*(l,r)
    =\gamma \max_{a'} q_*(h,a').
    \]
\end{solution}

\subsection{Упражнение 3.24(c. 96)}
Рисунок $3.5$ показывает, что оптимальная ценность лучшего состояния сеточного мира равна $24.4$ с точностью до одного знака после запятой. Воспользуйтесь своим знанием оптимальной стратегии и формулой $(3.8)$, чтобы выразить эту ценность символически, а затем вычислите ее с тремя знаками после запятой. 
\begin{figure}[H]
    \centering
    \includegraphics[width=0.75\linewidth]{fig_3_5.png}
    \caption{Рис. 3.5. Оптимальные решения задачи о сеточном мире}
\end{figure}
\begin{equation}
    G_t = R_{t+1} + \gamma R_{t+2} + \gamma^2 R_{t+3} + \cdots
= \sum_{k=0}^{\infty} \gamma^k R_{t+k+1}.
\tag{3.8}
\end{equation}
\begin{solution}
    Наилучшее состояние - $A$. Согласно оптимальной стратегии, из $A$
    агент получает вознаграждение $+10$ и мгновенно переходит в $A'$.
    Далее за $4$ шага с нулевыми наградами он возвращается в $A$,
    после чего цикл повторяется.
    
    Следовательно, начиная из $A$, ненулевые награды возникают
    через каждые $5$ шагов:
    \[
    10,\;\gamma^{5} \cdot 10,\;\gamma^{10}\cdot  10,\;\ldots
    \]
    
    По формуле (3.8)
    \[
    v_*(A)
    =
    \sum_{k=0}^{\infty} \gamma^k R_{t+k+1}
    =
    10 + \gamma^{5} \cdot 10 + \gamma^{10} \cdot 10 + \cdots
    =
    \frac{10}{1-\gamma^{5}}.
    \]
    
    При $\gamma = 0.9$ получаем
    \[
    v_*(A)
    =
    \frac{10}{1-0.9^{5}}
    =
    24.419.
    \]
\end{solution}

\subsection{Упражнение 3.25(c. 96)}
Запишите уравнение для $v_*$ в терминах $q_*$. 
\begin{solution}
    \[
    v_*(s)=\max_{a \in A(s)} q_*(s,a), \quad \forall s\in S.
    \]
\end{solution}

\subsection{Упражнение 3.26(c. 96)}
Запишите уравнение для $q_*$ в терминах $v_*$ и функции $p$ с четырьмя аргументами. 
\begin{solution}
    \[
    q_*(s,a)
    =
    \sum_{s',r}
    p(s',r \mid s,a)
    \left[
    r + \gamma v_*(s')
    \right],
    \quad \forall s \in S,\; a \in A(s).
    \]
\end{solution}

\subsection{Упражнение 3.27(c. 96)}
Запишите уравнение для $\pi_*$ в терминах $q_*$. 
\begin{solution}
    \[
    \pi_*(a\mid s)=
    \begin{cases}
    1, & a = \arg\max\limits_{a'} q_*(s,a'),\\[4pt]
    0, & \text{иначе},
    \end{cases}
    \qquad \forall s\in S,\; a\in A(s).
    \]
\end{solution}

\subsection{Упражнение 3.28(c. 96)}
Запишите уравнение для $\pi_*$ в терминах $v_*$ и функции $p$ с четырьмя аргументами.
\begin{solution}
    \[
    \begin{aligned}
    \pi_*(a \mid s)=
    &\begin{cases}
    1, 
    & a = \displaystyle
    \arg\max_{a'} \sum_{s',r} p(s',r \mid s,a')
    \bigl[ r + \gamma v_*(s') \bigr],
    \\[10pt]
    0,
    & \text{иначе},
    \end{cases} \\
    &\forall s \in S,\; a \in A(s).
    \end{aligned}
    \]
\end{solution}

\subsection{Упражнение 3.29(c. 96)}
Перепишите все четыре уравнения Беллмана для четырех функций ценности $(v_{\pi}, v_*, q_{\pi}, q_*)$ в  терминах функции $p$ с  тремя аргументами (3.4) и функции $r$ с двумя аргументами (3.5).
\begin{equation}
p(s' \mid s,a)
=
\Pr\{ S_t = s' \mid S_{t-1} = s,\; A_{t-1} = a \}
=
\sum_{r \in R} p(s',r \mid s,a).
\tag{3.4}
\end{equation}
\begin{equation}
r(s,a)
=
\mathbb{E}\!\left[ R_t \mid S_{t-1}=s,\; A_{t-1}=a \right]
=
\sum_{r \in R} r 
\sum_{s' \in S} 
p(s',r \mid s,a).
\tag{3.5}
\end{equation}
\begin{solution}
    \textbf{1. Уравнение Беллмана для $v_{\pi}$:}
    \[
    v_{\pi}(s)
    =
    \sum_{a} \pi(a \mid s)
    \left[
    r(s,a)
    +
    \gamma
    \sum_{s'} p(s' \mid s,a)\, v_{\pi}(s')
    \right].
    \]
    
    \textbf{2. Уравнение оптимальности для $v_*$:}
    \[
    v_{*}(s)
    =
    \max_{a}
    \left[
    r(s,a)
    +
    \gamma
    \sum_{s'} p(s' \mid s,a)\, v_{*}(s')
    \right].
    \]
    
    \textbf{3. Уравнение Беллмана для $q_{\pi}$:}
    \[
    q_{\pi}(s,a)
    =
    r(s,a)
    +
    \gamma
    \sum_{s'} p(s' \mid s,a)
    \sum_{a'} \pi(a' \mid s')\, q_{\pi}(s',a').
    \]
    
    \textbf{4. Уравнение оптимальности для $q_*$:}
    \[
    q_{*}(s,a)
    =
    r(s,a)
    +
    \gamma
    \sum_{s'} p(s' \mid s,a)
    \max_{a'} q_{*}(s',a').
    \]
\end{solution}
